% Unit 9 (Thermodynamics and Electrochemistry) guided notes -- 9 blocks.
\documentclass{apchem-notes}

\apcunit{9}
\apcunittitle{Thermodynamics and Electrochemistry}
\apcdoctitle{Guided Notes}
\apcsubtitle{CED 9.1--9.11 \textbullet\ Zumdahl Ch 17--18 \textbullet\ 9 blocks}

\begin{document}
\apctitleblock

\begin{apcnote}[One quantity, three languages]
Everything in this unit is a consequence of a single number, $\Delta G^\circ$,
and of the three equations that compute it or convert it:

\begin{center}
\renewcommand{\arraystretch}{1.25}
\begin{tabular}{@{}lll@{}}
\hline
\textbf{Equation} & \textbf{Answers the question} & \textbf{From} \\
\hline
$\Delta G^\circ = \Delta H^\circ - T\Delta S^\circ$ &
  is it favored, and at what $T$? & Unit 6 $+$ 9.1--9.2 \\
$\Delta G^\circ = -RT\ln K$ &
  how far does it go? & Unit 7 \\
$\Delta G^\circ = -nFE^\circ_{\text{cell}}$ &
  what voltage does it produce? & Unit 9 \\
\hline
\end{tabular}
\end{center}

Read across the middle column and you have the whole course: Unit 6 measured
heat, Unit 7 measured extent, Unit 5 measured rate. $\Delta G^\circ$ is what
connects the first two --- and topic 9.4 exists to insist that it says
nothing at all about the third.

\medskip
Three signs travel together and are worth memorizing as one fact:
\[ \Delta G^\circ < 0 \iff K > 1 \iff E^\circ_{\text{cell}} > 0 \]
\end{apcnote}

\begin{apctrap}[The word the CED retired]
The 2024 CED replaced ``spontaneous'' with
\textbf{thermodynamically favored} throughout, because ``spontaneous''
misleads students into hearing ``happens quickly, by itself.'' A favored
reaction may take geological time --- diamond turning into graphite is
favored and has not troubled anyone's jewellery.

Zumdahl's Chapter 17 is still titled ``Spontaneity, Entropy, and Free
Energy.'' Read the chapter, write the CED's word.
\end{apctrap}

% =====================================================================
\blocklesson{Entropy}{CED 9.1 \textbullet\ Zumdahl \S17.1}

\begin{objectives}
  \item State what entropy measures.
  \item Predict the sign of $\Delta S^\circ$ from the change described.
\end{objectives}

\reading{Zumdahl \S17.1}{835--841}

\segment{INSTRUCTION A}{25}{What entropy actually measures}

\notesection{Dispersal, not disorder}{9.1}
\practice{6}

\term{Entropy}, $S$, measures how many microscopic arrangements produce the
same macroscopic state --- how \answerblank[3cm]{dispersed} the energy and
matter of a system are. More available arrangements means
\answerblank[2.4cm]{higher} entropy.

\begin{apctrap}
``Disorder'' is the usual shorthand and it is a poor one. A messy room has
no more entropy than a tidy one, because nothing about the energy of its
molecules has changed. Argue from \textbf{how many ways the particles and
their energy can be arranged}, and from the volume available to them ---
that reasoning earns credit, and ``it looks messier'' does not.
\end{apctrap}

\notesection{The four situations to recognize}{9.1}

Nearly every $\Delta S$ question on the exam is one of these:

\begin{enumerate}[leftmargin=*,itemsep=0.35em]
  \item \textbf{Phase change.}
        solid $<$ liquid $\ll$ gas. Melting and boiling have
        $\Delta S$ \answerblank[2cm]{positive}; freezing and condensing,
        negative.
  \item \textbf{Change in moles of gas.} Count gas molecules on each side.
        More gas means $\Delta S$ \answerblank[2cm]{positive}. This
        dominates every other consideration --- if the mole count of gas
        changes, stop looking for other reasons.
  \item \textbf{Dissolving.} A crystal becoming dispersed ions is normally
        an \answerblank[2cm]{increase}.
  \item \textbf{Temperature.} Heating anything raises its entropy;
        molecules access more energy levels.
\end{enumerate}

\begin{apcnote}[The tie-breaker]
When two of these compete, \textbf{gas wins}. In
\ce{2H2(g) + O2(g) -> 2H2O(l)} three moles of gas become zero, so
$\Delta S^\circ$ is strongly negative even though a liquid is arguably
``more dispersed'' than a crystal would be.
\end{apcnote}

\segment{GUIDED PRACTICE}{15}{Predicting the sign}

Predict the sign of $\Delta S^\circ$ and give the controlling reason.

\begin{enumerate}[leftmargin=*,itemsep=0.4em]
  \item \ce{H2O(l) -> H2O(g)} \answerblank[6cm]{$+$; liquid becomes gas}
  \item \ce{2NO2(g) -> N2O4(g)} \answerblank[6cm]{$-$; 2 mol gas to 1}
  \item \ce{NaCl(s) -> Na+(aq) + Cl^-(aq)}
        \answerblank[6cm]{$+$; crystal disperses}
  \item \ce{CaCO3(s) -> CaO(s) + CO2(g)}
        \answerblank[6cm]{$+$; a gas appears}
  \item \ce{N2(g) + 3H2(g) -> 2NH3(g)}
        \answerblank[6cm]{$-$; 4 mol gas to 2}
\end{enumerate}

\segment{INSTRUCTION B}{20}{The second law}

\notesection{The universe, not the system}{9.1}
\practice{6}

The \term{second law of thermodynamics}: for any thermodynamically favored
process,
\[ \Delta S_{\text{univ}} = \Delta S_{\text{sys}} + \Delta S_{\text{surr}}
   > 0 \]

The system's entropy is allowed to \emph{decrease} --- water freezes, ice
crystals are more ordered than liquid --- provided the surroundings gain
more than the system loses. Freezing releases heat to the surroundings, and
that heat raises their entropy.

\begin{apcnote}[Why this becomes Gibbs free energy]
Tracking the surroundings is a nuisance. The next block folds them into a
single system-only quantity, which is the entire reason $\Delta G$ is worth
defining. If you understand $\Delta S_{\text{univ}} > 0$, you already
understand $\Delta G^\circ < 0$ --- they are the same statement multiplied
by $-T$.
\end{apcnote}

\segment{APPLICATION}{20}{Reasoning with the second law}

\begin{enumerate}[leftmargin=*,itemsep=0.4em]
  \item Water freezing at \SI{-10}{\celsius} is thermodynamically favored,
        yet the water's entropy decreases. Explain how both are true.
        \answerlines[3]{Freezing releases heat to the surroundings. Below
        \SI{0}{\celsius} the entropy gained by the surroundings from that
        released heat exceeds the entropy lost by the water, so
        $\Delta S_{\text{univ}}$ is still positive and the second law is
        satisfied.}
  \item A student says a reaction with $\Delta S_{\text{sys}} < 0$ can
        never happen. Correct them in one sentence.
        \answerlines[2]{The second law constrains the entropy of the
        \emph{universe}, not of the system; a system may lose entropy as
        long as the surroundings gain more.}
  \item Rank by absolute entropy at the same temperature: \ce{H2O(s)},
        \ce{H2O(l)}, \ce{H2O(g)}. \answerblank[7cm]{\ce{H2O(s)} $<$
        \ce{H2O(l)} $<$ \ce{H2O(g)}}
\end{enumerate}

\begin{exitticket}
Predict the sign of $\Delta S^\circ$ for
\ce{2SO2(g) + O2(g) -> 2SO3(g)} and justify in one sentence.
\keyonly{\begin{apcanswer}Negative. Three moles of gas become two, and a
decrease in the number of gas particles dominates any other
contribution.\end{apcanswer}}
\end{exitticket}

% =====================================================================
\blocklesson{Absolute Entropy and $\Delta S^\circ$}{CED 9.2 \textbullet\ Zumdahl \S17.2, \S17.5--17.6}

\begin{objectives}
  \item Explain why entropy has an absolute scale but enthalpy does not.
  \item Calculate $\Delta S^\circ$ from tabulated absolute entropies.
\end{objectives}

\reading{Zumdahl \S17.2, \S17.5--17.6}{841--842, 848--853}

\begin{warmup}[4]
  \item Sign of $\Delta S^\circ$ for \ce{2NO2(g) -> N2O4(g)}:
        \answerblank[2cm]{negative}
  \item Which has more entropy, \ce{H2O(l)} or \ce{H2O(g)}?
        \answerblank[2.4cm]{\ce{H2O(g)}}
  \item $\Delta S_{\text{univ}}$ for a favored process:
        \answerblank[2cm]{$> 0$}
  \item The AP term replacing ``spontaneous'':
        \answerblank[5.4cm]{thermodynamically favored}
\end{warmup}

\segment{INSTRUCTION A}{25}{A real zero}

\notesection{The third law}{9.2}
\practice{5}

The \term{third law of thermodynamics} says a perfect crystal at
\SI{0}{\kelvin} has entropy exactly \answerblank[1.4cm]{zero}. That gives
entropy a genuine origin, so every substance has a tabulated
\term{absolute entropy} $S^\circ$ --- always a positive number.

\begin{apctrap}
Contrast with enthalpy. There is no absolute $H$, so tables list
$\Delta H_f^\circ$ and assign zero to elements in their standard states.

Entropy tables are different: \textbf{$S^\circ$ of an element is not zero.}
\ce{O2(g)} has $S^\circ = \SI{205}{\joule\per\kelvin\per\mole}$. Setting
element entropies to zero out of habit is a reliable way to lose a point.
\end{apctrap}

\[ \boxed{\;\Delta S^\circ_{\text{rxn}} = \sum n S^\circ_{\text{products}}
   - \sum n S^\circ_{\text{reactants}}\;} \]

Coefficients multiply, exactly as in Hess's law. And note the units:
$S^\circ$ is in \answerblank[2cm]{joules} while $\Delta H^\circ$ is in
\answerblank[2.4cm]{kilojoules}.

\begin{workedexample}[Worked example 1: decomposing limestone]
\[ \ce{CaCO3(s) -> CaO(s) + CO2(g)} \]
$S^\circ$: \ce{CaCO3} 93, \ce{CaO} 40, \ce{CO2} 214
\si{\joule\per\kelvin\per\mole}.
\[ \Delta S^\circ = (40 + 214) - 93 = \mathbf{+161}~\si{\joule\per\kelvin} \]
Positive, as predicted --- a gas appeared where there was none.
\end{workedexample}

\begin{workedexample}[Worked example 2: the Haber process]
\[ \ce{N2(g) + 3H2(g) -> 2NH3(g)} \]
$S^\circ$: \ce{N2} 192, \ce{H2} 131, \ce{NH3} 193.
\[ \Delta S^\circ = 2(193) - [192 + 3(131)] = 386 - 585
   = \mathbf{-199}~\si{\joule\per\kelvin} \]
Four moles of gas became two, so the sign had to be negative. Dropping the
coefficient 3 on hydrogen gives $-68$ --- wrong magnitude entirely.
\end{workedexample}

\segment{GUIDED PRACTICE}{15}{Running the sums}

$S^\circ$ (\si{\joule\per\kelvin\per\mole}): \ce{H2} 131, \ce{O2} 205,
\ce{H2O(l)} 70, \ce{H2O(g)} 189, \ce{NO2} 240, \ce{N2O4} 304.

\begin{enumerate}[leftmargin=*,itemsep=0.4em]
  \item \ce{2H2(g) + O2(g) -> 2H2O(l)}: \workspace[2cm]
        \keyonly{\begin{apcanswer}$2(70) - [2(131) + 205] = 140 - 467
        = \mathbf{-327}~\si{\joule\per\kelvin}$\end{apcanswer}}
  \item \ce{H2O(l) -> H2O(g)}: \answerblank[3.4cm]{$+119$ J/K}
  \item \ce{2NO2(g) -> N2O4(g)}: \answerblank[3.4cm]{$-176$ J/K}
\end{enumerate}

\segment{APPLICATION}{20}{Interpreting the numbers}

\begin{enumerate}[leftmargin=*,itemsep=0.4em]
  \item Why is $S^\circ$ for \ce{H2O(g)} nearly triple that of
        \ce{H2O(l)}, when it is the same substance?
        \answerlines[2]{In the gas the molecules occupy a far larger volume
        and move independently, so vastly more microscopic arrangements
        correspond to the same macroscopic state.}
  \item A student computes the Haber $\Delta S^\circ$ as
        $193 - (192 + 131) = -130$~J/K. Find both errors.
        \answerlines[3]{They omitted the coefficient 2 on \ce{NH3} and the
        coefficient 3 on \ce{H2}. Correctly,
        $2(193) - [192 + 3(131)] = -199$~J/K.}
  \item For \ce{H2O(l) -> H2O(g)}, $\Delta S^\circ = +119$~J/K and
        $\Delta H^\circ = +44$~kJ. Say what each sign contributes to
        favorability at \SI{25}{\celsius}.
        \answerlines[3]{Entropy favors vaporization --- gas is far more
        dispersed. Enthalpy opposes it, since hydrogen bonds must break. At
        \SI{25}{\celsius} the enthalpy term wins, so water does not all
        become vapour; raising $T$ grows the $T\Delta S$ term until it
        does.}
\end{enumerate}

\begin{exitticket}
Why can a table list an absolute $S^\circ$ for oxygen gas but not an
absolute $H^\circ$?
\keyonly{\begin{apcanswer}The third law fixes a true zero for entropy --- a
perfect crystal at \SI{0}{\kelvin} has $S = 0$ --- so entropies are
measurable on an absolute scale. Enthalpy has no such natural zero, so only
\emph{changes} are measurable and tables use $\Delta H_f^\circ$ with
elements assigned zero by convention.\end{apcanswer}}
\end{exitticket}

% =====================================================================
\blocklesson{Gibbs Free Energy}{CED 9.3 \textbullet\ Zumdahl \S17.3--17.4}

\begin{objectives}
  \item Calculate $\Delta G^\circ$ from $\Delta H^\circ$ and
        $\Delta S^\circ$.
  \item Determine the temperature range over which a process is favored.
\end{objectives}

\reading{Zumdahl \S17.3--17.4}{842--848}

\begin{warmup}[4]
  \item $\Delta S^\circ$ for \ce{CaCO3(s) -> CaO(s) + CO2(g)}:
        \answerblank[2.6cm]{$+161$ J/K}
  \item $\Delta S_{\text{univ}}$ must be what sign for a favored process?
        \answerblank[2cm]{positive}
  \item Units of $S^\circ$ versus $\Delta H^\circ$:
        \answerblank[5.4cm]{J/(K$\cdot$mol) versus kJ/mol}
  \item $S^\circ$ of an element in its standard state is:
        \answerblank[3.4cm]{not zero}
\end{warmup}

\segment{INSTRUCTION A}{25}{One number that decides}

\notesection{Folding the surroundings in}{9.3}
\practice{5}

\term{Gibbs free energy} removes the need to track the surroundings:
\[ G = H - TS \qquad\Longrightarrow\qquad
   \boxed{\;\Delta G^\circ = \Delta H^\circ - T\Delta S^\circ\;} \]

Dividing $\Delta G = \Delta H - T\Delta S$ by $-T$ returns exactly
$\Delta S_{\text{surr}} + \Delta S_{\text{sys}} = \Delta S_{\text{univ}}$.
So $\Delta G$ is the second law rewritten \emph{entirely in terms of the
system}.

\begin{center}
\fbox{\parbox{0.8\linewidth}{\centering
$\Delta G^\circ < 0$: \textbf{thermodynamically favored} \\
$\Delta G^\circ > 0$: not favored (the reverse is) \\
$\Delta G^\circ = 0$: at equilibrium}}
\end{center}

The sign flip is worth stating aloud: $\Delta S_{\text{univ}}$ must be
\answerblank[1.8cm]{positive}, but $\Delta G^\circ$ must be
\answerblank[1.8cm]{negative}, because of the $-T$ used in the derivation.

\begin{apctrap}
\textbf{Convert the entropy to kJ before subtracting.} $\Delta H^\circ$ is
in kJ and $\Delta S^\circ$ in J/K. Forgetting the factor of 1000 is the
single most common arithmetic error in this unit, and it produces answers
wrong by roughly a factor of a thousand --- large enough to spot if you
sanity-check the magnitude.
\end{apctrap}

\segment{INSTRUCTION B}{20}{When temperature decides}

\notesection{The four cases}{9.3}
\practice{5}

Because $\Delta G^\circ = \Delta H^\circ - T\Delta S^\circ$, the two signs
between them decide whether temperature matters at all:

\begin{center}
\renewcommand{\arraystretch}{1.2}
\begin{tabular}{@{}cccl@{}}
\hline
$\Delta H^\circ$ & $\Delta S^\circ$ & \textbf{Favored at} &
  \textbf{Comment} \\
\hline
$-$ & $+$ & \textbf{all temperatures} & both terms help; no calculation
  needed \\
$+$ & $-$ & \textbf{no temperature} & both terms oppose; never favored \\
$+$ & $+$ & \textbf{high $T$} & entropy wins once $T$ is large enough \\
$-$ & $-$ & \textbf{low $T$} & enthalpy wins while $T$ is small \\
\hline
\end{tabular}
\end{center}

For the two mixed cases the \term{crossover temperature}, where
$\Delta G^\circ = 0$, is
\[ T = \frac{\Delta H^\circ}{\Delta S^\circ} \]

\begin{workedexample}[Worked example 3: why limestone needs a kiln]
\ce{CaCO3(s) -> CaO(s) + CO2(g)}: $\Delta H^\circ = +178.5$~kJ,
$\Delta S^\circ = +161$~J/K. Both positive $\Rightarrow$ favored only at
high temperature.

\textbf{At \SI{298}{\kelvin}:}
$\Delta G^\circ = 178.5 - 298(0.161) = +131$~kJ --- not favored, which is
why limestone buildings do not crumble into quicklime.

\textbf{Crossover:}
\[ T = \frac{178.5}{0.161} = \mathbf{1109}~\si{\kelvin}
   \quad (\approx \SI{836}{\celsius}) \]
Which is about the temperature at which lime kilns are operated.
\end{workedexample}

\segment{APPLICATION}{20}{Predicting and calculating}

\begin{enumerate}[leftmargin=*,itemsep=0.4em]
  \item For the Haber process $\Delta H^\circ = -92$~kJ and
        $\Delta S^\circ = -199$~J/K. Find $\Delta G^\circ$ at
        \SI{298}{\kelvin} and the crossover temperature. \workspace[2.8cm]
        \keyonly{\begin{apcanswer}$\Delta G^\circ = -92 - 298(-0.199)
        = -92 + 59 = \mathbf{-33}$~kJ, so favored at \SI{298}{\kelvin}.
        Crossover $T = 92/0.199 = \mathbf{462}$~K --- favored \emph{below}
        462 K, since both terms are negative.\end{apcanswer}}
  \item For \ce{H2O(l) -> H2O(g)}, $\Delta H^\circ = +44$~kJ and
        $\Delta S^\circ = +119$~J/K. Find the crossover temperature and say
        what it physically is. \workspace[2.4cm]
        \keyonly{\begin{apcanswer}$T = 44/0.119 = \mathbf{370}$~K
        $\approx \SI{97}{\celsius}$ --- the \textbf{boiling point}. Above
        it vaporization is favored at 1 atm; below it condensation is. The
        tabulated value is 373 K; the gap is rounding in the
        table.\end{apcanswer}}
  \item Without calculating, say whether
        \ce{2H2O2(l) -> 2H2O(l) + O2(g)} ($\Delta H^\circ < 0$) is favored,
        and at what temperatures. \answerlines[2]{$\Delta H^\circ$ is
        negative and $\Delta S^\circ$ positive (a gas forms from a liquid),
        so both terms help and it is favored at \emph{all} temperatures. No
        calculation is needed.}
\end{enumerate}

\begin{exitticket}
A reaction has $\Delta H^\circ = +50$~kJ and $\Delta S^\circ = +100$~J/K.
Is it favored at \SI{298}{\kelvin}? Above what temperature does that change?
\keyonly{\begin{apcanswer}$\Delta G^\circ = 50 - 298(0.100) = +20$~kJ, so
not favored. Crossover at $T = 50/0.100 = \SI{500}{\kelvin}$; above
\SI{500}{\kelvin} it becomes favored.\end{apcanswer}}
\end{exitticket}

% =====================================================================
\blocklesson{Formation Free Energies and Kinetic Control}{CED 9.3, 9.4 \textbullet\ Zumdahl \S17.4, \S17.7}

\begin{objectives}
  \item Calculate $\Delta G^\circ$ from free energies of formation.
  \item Explain why a favored reaction may proceed immeasurably slowly.
\end{objectives}

\reading{Zumdahl \S17.4, \S17.7}{845--848, 853--859}

\begin{warmup}[4]
  \item $\Delta G^\circ = $ \answerblank[4cm]{$\Delta H^\circ -
        T\Delta S^\circ$}
  \item Favored means $\Delta G^\circ$ is: \answerblank[2cm]{negative}
  \item $\Delta H^\circ < 0$ and $\Delta S^\circ > 0$: favored at
        \answerblank[4cm]{all temperatures}
  \item Crossover temperature $= $ \answerblank[3cm]{$\Delta H^\circ /
        \Delta S^\circ$}
\end{warmup}

\segment{INSTRUCTION A}{25}{The second route to $\Delta G^\circ$}

\notesection{Free energies of formation}{9.3}
\practice{5}

Exactly as with enthalpy, $\Delta G^\circ$ can be assembled from tabulated
formation values:
\[ \Delta G^\circ_{\text{rxn}} = \sum n\Delta G_f^\circ{}_{\text{products}}
   - \sum n\Delta G_f^\circ{}_{\text{reactants}} \]
and $\Delta G_f^\circ$ of an element in its standard state is
\answerblank[1.4cm]{zero} --- unlike $S^\circ$.

\begin{workedexample}[Worked example 4: two routes, one answer]
For \ce{2H2(g) + O2(g) -> 2H2O(l)}:

\textbf{Route 1 --- from $\Delta H^\circ$ and $\Delta S^\circ$.}
$\Delta H^\circ = 2(-286) = -572$~kJ;
$\Delta S^\circ = -327$~J/K $= -0.327$~kJ/K.
\[ \Delta G^\circ = -572 - 298(-0.327) = -572 + 97
   = \mathbf{-475}~\text{kJ} \]

\textbf{Route 2 --- from $\Delta G_f^\circ$.}
$\Delta G^\circ = 2(-237) - 0 = \mathbf{-474}$~kJ.

The two agree to within table rounding. Note again that the entropy had to
be converted from J/K to kJ/K first.
\end{workedexample}

\segment{INSTRUCTION B}{20}{Favored is not the same as fast}

\notesection{Thermodynamic versus kinetic control}{9.4}
\practice{6}

This is the point of topic 9.4, and it is the one place where this unit and
Unit 5 must be held apart:

\begin{center}
\fbox{\parbox{0.9\linewidth}{\centering
$\Delta G^\circ$ tells you \textbf{whether} and \textbf{how far}. \\
$E_a$ tells you \textbf{how fast}. \\
They are independent.}}
\end{center}

A reaction that is thermodynamically favored but has a large
\answerblank[3.6cm]{activation energy} proceeds too slowly to observe. It is
said to be under \term{kinetic control}.

\begin{itemize}[leftmargin=*,itemsep=0.3em]
  \item \textbf{Diamond to graphite.} $\Delta G^\circ$ is negative, so
        graphite is the favored form at room conditions. The rearrangement
        requires breaking a covalent network, so it does not measurably
        occur.
  \item \textbf{Petrol in air.} Combustion is strongly favored; a tank of
        fuel is stable because the reaction needs a spark to clear
        $E_a$.
\end{itemize}

\begin{apctrap}
A \textbf{catalyst} lowers $E_a$ and speeds a reaction up. It does
\emph{not} change $\Delta G^\circ$, $\Delta H^\circ$, $\Delta S^\circ$, or
\answerblank[1.4cm]{$K$}. A catalyst cannot make an unfavorable reaction
favorable --- it can only let an already-favored one arrive sooner.
\end{apctrap}

\segment{APPLICATION}{20}{Thermodynamics versus kinetics}

\begin{enumerate}[leftmargin=*,itemsep=0.4em]
  \item A reaction has $\Delta G^\circ = -150$~kJ but no measurable rate at
        room temperature. Is this a contradiction? Explain.
        \answerlines[3]{No. $\Delta G^\circ$ describes the relative free
        energies of reactants and products and says nothing about the
        barrier between them. A large activation energy makes the reaction
        slow while leaving it thermodynamically favored --- the reaction is
        under kinetic control.}
  \item Give two ways to speed that reaction up, and state the effect of
        each on $K$. \answerlines[3]{Raise the temperature --- more
        molecules exceed $E_a$; this does change $K$, since $K$ is
        temperature-dependent. Add a catalyst --- this lowers $E_a$ and
        leaves $K$ unchanged.}
  \item Explain why ``spontaneous'' is a poor word for $\Delta G^\circ < 0$
        using diamond as the example.
        \answerlines[2]{It suggests the reaction happens by itself and
        quickly. Diamond converting to graphite is thermodynamically
        favored yet immeasurably slow, so ``thermodynamically favored'' is
        the accurate description.}
\end{enumerate}

\begin{exitticket}
State one thing a catalyst changes and two things it does not.
\keyonly{\begin{apcanswer}It lowers the activation energy (and so raises
the rate). It does not change $\Delta G^\circ$ or $K$ --- nor
$\Delta H^\circ$ or $\Delta S^\circ$.\end{apcanswer}}
\end{exitticket}

% =====================================================================
\blocklesson{Free Energy and Equilibrium}{CED 9.5 \textbullet\ Zumdahl \S17.9}

\begin{objectives}
  \item Relate $\Delta G^\circ$ to $K$ and estimate one from the other.
  \item Distinguish $\Delta G$ from $\Delta G^\circ$ using $Q$.
\end{objectives}

\reading{Zumdahl \S17.9}{862--867}

\begin{warmup}[4]
  \item If $K > 1$, which side is favored? \answerblank[2.4cm]{products}
  \item $Q$ is calculated like $K$ but with:
        \answerblank[6.3cm]{non-equilibrium concentrations}
  \item What raises a rate without changing $K$?
        \answerblank[2.6cm]{a catalyst}
  \item $\Delta G^\circ$ for a favored process is:
        \answerblank[2cm]{negative}
\end{warmup}

\segment{INSTRUCTION A}{25}{Two languages for one fact}

\notesection{$\Delta G^\circ$ and $K$}{9.5}
\practice{5}

You now have two ways of saying ``products are favored'': $K > 1$ from
Unit 7, and $\Delta G^\circ < 0$ from this unit. They are the same
statement:
\[ \boxed{\;\Delta G^\circ = -RT\ln K\;}
   \qquad\text{equivalently}\qquad K = e^{-\Delta G^\circ/RT} \]
with $R = \SI{8.314}{\joule\per\mole\per\kelvin}$ --- note
\answerblank[1.8cm]{joules}, so $\Delta G^\circ$ must be converted from kJ.

\begin{center}
\renewcommand{\arraystretch}{1.2}
\begin{tabular}{@{}ccl@{}}
\hline
$\Delta G^\circ$ & $K$ & \\
\hline
negative & $> 1$ & products favored at equilibrium \\
zero     & $= 1$ & neither side favored \\
positive & $< 1$ & reactants favored \\
\hline
\end{tabular}
\end{center}

\begin{apcnote}[The MCQ half is no-calculator]
You are not asked to evaluate logarithms by hand. You are asked to know
which \emph{way} the relationship runs: more negative $\Delta G^\circ$ means
larger $K$, and the dependence is exponential --- so a modest change in
$\Delta G^\circ$ produces an enormous change in $K$. Reason with the sign
and the direction, not with arithmetic.
\end{apcnote}

\segment{INSTRUCTION B}{20}{$\Delta G$ versus $\Delta G^\circ$}

\notesection{Away from standard conditions}{9.5}
\practice{6}

$\Delta G^\circ$ is a fixed property of the reaction at a given temperature
--- every species at standard conditions. $\Delta G$ is what the reaction
mixture \emph{actually} has right now, and depends on the current
composition through $Q$:
\[ \Delta G = \Delta G^\circ + RT\ln Q \]

\begin{itemize}[leftmargin=*,itemsep=0.3em]
  \item $Q < K$: the mixture shifts \answerblank[2.4cm]{forward} and
        $\Delta G < 0$.
  \item $Q = K$: \answerblank[3cm]{equilibrium}, and $\Delta G = 0$.
  \item $Q > K$: the mixture shifts \answerblank[2.4cm]{reverse} and
        $\Delta G > 0$.
\end{itemize}

\begin{apctrap}
At equilibrium $\Delta G = 0$, \textbf{not} $\Delta G^\circ$. $\Delta G^\circ$
is zero only in the special case $K = 1$. Writing ``$\Delta G^\circ = 0$ at
equilibrium'' is a standard lost point --- and setting $\Delta G = 0$ in
$\Delta G = \Delta G^\circ + RT\ln Q$ is precisely how
$\Delta G^\circ = -RT\ln K$ is derived.
\end{apctrap}

\segment{APPLICATION}{20}{Connecting the two}

\begin{enumerate}[leftmargin=*,itemsep=0.4em]
  \item Reaction A has $\Delta G^\circ = -5$~kJ, reaction B has
        $\Delta G^\circ = -50$~kJ. Compare their $K$ values qualitatively.
        \answerlines[3]{Both have $K > 1$. B's is very much larger --- the
        relationship is exponential, so a tenfold more negative
        $\Delta G^\circ$ produces a vastly larger $K$, not a tenfold larger
        one. A's is only modestly above 1.}
  \item A reaction has $\Delta G^\circ = +20$~kJ. Can it ever proceed
        forward? \answerlines[3]{Yes. $\Delta G^\circ$ refers to standard
        conditions. If $Q$ is small enough --- very little product present
        --- then $RT\ln Q$ is sufficiently negative to make
        $\Delta G < 0$, and the reaction proceeds until $Q$ reaches $K$.}
  \item At equilibrium, state the values of $\Delta G$ and of $Q$.
        \answerblank[6cm]{$\Delta G = 0$ and $Q = K$}
\end{enumerate}

\begin{exitticket}
A reaction has $K = 1$. What is $\Delta G^\circ$, and is the reaction
favored?
\keyonly{\begin{apcanswer}$\Delta G^\circ = -RT\ln(1) = 0$. Neither
direction is favored --- at standard conditions the reaction is already at
equilibrium.\end{apcanswer}}
\end{exitticket}

% =====================================================================
\blocklesson{Dissolution and Coupled Reactions}{CED 9.6, 9.7 \textbullet\ Zumdahl \S17.5, \S17.10}

\begin{objectives}
  \item Identify the enthalpy and entropy contributions to dissolution.
  \item Explain how coupling or external energy drives an unfavorable
        process.
\end{objectives}

\reading{Zumdahl \S17.5, \S17.10}{848--849, 867--869}

\begin{warmup}[4]
  \item $\Delta G^\circ = -RT\ln K$; $R$ is in:
        \answerblank[2.4cm]{J/(mol$\cdot$K)}
  \item At equilibrium $\Delta G = $ \answerblank[1.4cm]{0}
  \item $\Delta S$ for dissolving a crystal is usually:
        \answerblank[2cm]{positive}
  \item A catalyst changes: \answerblank[3cm]{$E_a$ only}
\end{warmup}

\segment{INSTRUCTION A}{25}{Why some salts dissolve}

\notesection{Three competing contributions}{9.6}
\practice{4}

Dissolving a salt is a free-energy balance with three parts:

\begin{enumerate}[leftmargin=*,itemsep=0.3em]
  \item \textbf{Breaking up the lattice.} $\Delta H$ strongly
        \answerblank[2cm]{positive}, $\Delta S$ positive.
  \item \textbf{Reorganizing the solvent} around the ions. $\Delta S$
        \answerblank[2cm]{negative} --- water becomes ordered in hydration
        shells.
  \item \textbf{Ion--solvent attraction} (hydration). $\Delta H$
        \answerblank[2cm]{negative}.
\end{enumerate}

\begin{apcnote}[Why the CED does not ask you to predict the total]
The framework says outright that predicting the overall $\Delta G^\circ$ of
dissolution is challenging because of the cancellations among these three
factors. The enthalpy terms are large and opposite; so are the entropy
terms. What remains is a small difference between big numbers.

You are expected to \textbf{identify the contributions and their signs},
not to predict solubility from scratch. Reasoning through the three factors
earns credit even where the net direction is uncertain.
\end{apcnote}

\begin{workedexample}[Worked example 5: dissolving that cools the beaker]
\ce{NH4NO3} dissolves \emph{endothermically} --- the beaker gets cold, so
$\Delta H^\circ > 0$. Yet it is very soluble. Why?

Because $\Delta S^\circ$ is strongly positive: an ordered crystal becomes
freely moving hydrated ions dispersed through the solvent. At room
temperature $T\Delta S^\circ$ exceeds $\Delta H^\circ$, so
$\Delta G^\circ = \Delta H^\circ - T\Delta S^\circ < 0$.

\textbf{Entropy is driving this one alone}, with the enthalpy term working
against it --- the case to remember whenever you are tempted to equate
``exothermic'' with ``favored''.
\end{workedexample}

\segment{INSTRUCTION B}{20}{Driving the unfavorable}

\notesection{Coupling and external energy}{9.7}
\practice{4}

A process with $\Delta G^\circ > 0$ will not happen on its own. There are
two ways to make it happen anyway.

\textbf{1. Supply external energy.}
\begin{itemize}[leftmargin=*,itemsep=0.25em]
  \item \textbf{Electrical energy} driving an electrolytic cell or charging
        a battery --- the subject of the next three blocks.
  \item \textbf{Light} driving photosynthesis, which builds glucose from
        \ce{CO2} and water with a hugely positive $\Delta G^\circ$.
\end{itemize}

\textbf{2. Couple it to a favorable reaction.} Because free energies
\answerblank[2cm]{add}, an unfavorable step can be carried by a favorable
one --- provided the two share a \answerblank[4.6cm]{common intermediate}
and the \emph{sum} has $\Delta G^\circ < 0$.

\begin{workedexample}[Worked example 6: how cells pay for chemistry]
Phosphorylating glucose is unfavorable on its own:
\[ \text{glucose} + \ce{P_i} \to \text{glucose-6-phosphate}
   \qquad \Delta G^\circ = +13.8~\text{kJ/mol} \]
Hydrolysing ATP is strongly favorable:
\[ \ce{ATP} + \ce{H2O} \to \ce{ADP} + \ce{P_i}
   \qquad \Delta G^\circ = -30.5~\text{kJ/mol} \]
Coupled through the shared \ce{P_i}:
\[ \Delta G^\circ = +13.8 + (-30.5) = \mathbf{-16.7}~\text{kJ/mol} \]
and therefore favored. This is what ``ATP is the cell's energy currency''
means --- not that ATP contains energy, but that its hydrolysis is
favorable enough to pay for reactions that are not.
\end{workedexample}

\segment{APPLICATION}{20}{Putting free energy to work}

\begin{enumerate}[leftmargin=*,itemsep=0.4em]
  \item \ce{NaCl} dissolves with $\Delta H^\circ$ very slightly positive
        and still dissolves readily. Explain using the three contributions.
        \answerlines[3]{The lattice energy absorbed and the hydration
        energy released nearly cancel, leaving $\Delta H^\circ$ near zero.
        The entropy increase from an ordered crystal becoming dispersed
        hydrated ions is large enough that $-T\Delta S^\circ$ makes
        $\Delta G^\circ$ negative.}
  \item Reaction X has $\Delta G^\circ = +25$~kJ. It is coupled to
        reaction Y with $\Delta G^\circ = -40$~kJ. Is the coupled process
        favored, and what must the two share? \workspace[2cm]
        \keyonly{\begin{apcanswer}$\Delta G^\circ = +25 + (-40) =
        \mathbf{-15}$~kJ, so yes. They must share a \textbf{common
        intermediate} --- a species produced by one and consumed by the
        other --- or the coupling is only arithmetic, not
        chemical.\end{apcanswer}}
  \item Name the external energy source that drives each: electrolysis of
        water; photosynthesis. \answerblank[6.2cm]{electrical; light}
\end{enumerate}

\begin{exitticket}
State the one condition that makes coupling chemically real rather than
just an addition of two numbers.
\keyonly{\begin{apcanswer}The two reactions must share a \textbf{common
intermediate}, so that the product of one is consumed by the other in the
same system. Without it the free energies still add on paper, but nothing
connects the two processes.\end{apcanswer}}
\end{exitticket}

% =====================================================================
\blocklesson{Galvanic and Electrolytic Cells}{CED 9.8 \textbullet\ Zumdahl \S18.1, \S18.7}

\begin{objectives}
  \item Identify anode, cathode, and the direction of electron flow.
  \item Distinguish galvanic from electrolytic cells.
\end{objectives}

\reading{Zumdahl \S18.1, \S18.7}{882--884, 906--911}

\begin{warmup}[4]
  \item Oxidation is loss of: \answerblank[2cm]{electrons}
  \item An unfavorable process has $\Delta G^\circ$:
        \answerblank[2cm]{positive}
  \item One way to drive an unfavorable reaction:
        \answerblank[4.7cm]{supply external energy}
  \item Oxidation number of \ce{Zn} in \ce{Zn^2+}:
        \answerblank[1.4cm]{$+2$}
\end{warmup}

\segment{INSTRUCTION A}{25}{Separating the half-reactions}

\notesection{Why a wire is involved at all}{9.8}
\practice{1}

Drop zinc into copper sulfate and electrons pass directly from zinc to
\ce{Cu^2+}; the energy is released as \answerblank[1.6cm]{heat}. Separate
the two half-reactions into different compartments and the electrons must
travel through a \answerblank[1.6cm]{wire} instead --- and along the way
they can do electrical work. That is the whole idea of a cell.

\notesection{The parts, and what each does}{9.8}

\begin{center}
\renewcommand{\arraystretch}{1.25}
\begin{tabular}{@{}p{2.6cm}p{4.2cm}p{7.4cm}@{}}
\hline
\textbf{Part} & \textbf{Where} & \textbf{What happens} \\
\hline
Anode & one half-cell &
  \answerblank[3.2cm]{oxidation}; electrons leave \\
Cathode & the other half-cell &
  \answerblank[3.2cm]{reduction}; electrons arrive \\
Wire & between electrodes & carries electrons, anode $\to$ cathode \\
Salt bridge & joins the solutions &
  carries ions to keep both solutions neutral \\
\hline
\end{tabular}
\end{center}

\textbf{An~Ox and Red~Cat} --- \textbf{An}ode \textbf{Ox}idation,
\textbf{Red}uction \textbf{Cat}hode --- holds in every cell, galvanic or
electrolytic. It is the one fact in this block that never changes.

\begin{apcnote}[What the CED excludes here]
Topic 9.8 \textbf{excludes} assigning positive and negative labels to the
electrodes. The signs actually reverse between galvanic and electrolytic
cells, which is exactly why the framework leaves them out. Identify
electrodes by \emph{anode} and \emph{cathode} and by what happens there;
never by charge.
\end{apcnote}

Without the salt bridge, oxidation at the anode would leave that solution
building up \answerblank[3cm]{positive charge} and the cathode solution
going negative. Charge separation halts electron flow almost immediately,
so the cell stops.

\segment{INSTRUCTION B}{20}{The same hardware, run backwards}

\notesection{Galvanic against electrolytic}{9.8}
\practice{1}

\begin{center}
\renewcommand{\arraystretch}{1.25}
\begin{tabular}{@{}p{3.2cm}p{5.2cm}p{5.2cm}@{}}
\hline
 & \textbf{Galvanic (voltaic)} & \textbf{Electrolytic} \\
\hline
$\Delta G^\circ$ & \answerblank[2.4cm]{negative} &
  \answerblank[2.4cm]{positive} \\
$E^\circ_{\text{cell}}$ & \answerblank[2.4cm]{positive} &
  \answerblank[2.4cm]{negative} \\
Drives itself? & yes --- it is a battery & no --- needs a power supply \\
Energy & chemical $\to$ electrical & electrical $\to$ chemical \\
Example & a AA cell discharging & electroplating; recharging \\
\hline
\end{tabular}
\end{center}

An electrolytic cell is CED topic 9.7's ``external energy source driving an
unfavorable process,'' built in hardware.

\segment{APPLICATION}{20}{Describing cells}

\begin{enumerate}[leftmargin=*,itemsep=0.4em]
  \item In a galvanic cell made from \ce{Zn}/\ce{Zn^2+} and
        \ce{Cu}/\ce{Cu^2+}, zinc is oxidized. Name the anode, the cathode,
        and the direction of electron flow.
        \answerlines[2]{Zinc is the anode (oxidation); copper is the
        cathode (reduction). Electrons flow through the wire from the zinc
        electrode to the copper electrode.}
  \item State what happens to the mass of each electrode as that cell runs.
        \answerlines[2]{The zinc anode loses mass as it is oxidized to
        \ce{Zn^2+}. The copper cathode gains mass as \ce{Cu^2+} is reduced
        and plates onto it.}
  \item A student labels the anode ``negative'' on an exam question about
        an electrolytic cell. Explain why the CED avoids this.
        \answerlines[3]{The sign convention reverses between cell types:
        the anode is negative in a galvanic cell but positive in an
        electrolytic one, because in the latter it is connected to the
        positive terminal of the supply. Since oxidation occurs at the
        anode in both, identifying electrodes by process rather than charge
        is unambiguous.}
\end{enumerate}

\begin{exitticket}
Give the one statement about anode and cathode that is true in every cell,
and one statement that is not.
\keyonly{\begin{apcanswer}Always true: oxidation occurs at the anode and
reduction at the cathode. Not always true: that the anode is the negative
electrode --- that holds for galvanic cells and reverses for electrolytic
ones.\end{apcanswer}}
\end{exitticket}

% =====================================================================
\blocklesson{Cell Potential and Free Energy}{CED 9.9 \textbullet\ Zumdahl \S18.2--18.3}

\begin{objectives}
  \item Calculate $E^\circ_{\text{cell}}$ from standard reduction
        potentials.
  \item Relate $E^\circ_{\text{cell}}$ to $\Delta G^\circ$ and to $K$.
\end{objectives}

\reading{Zumdahl \S18.2--18.3}{884--895}

\begin{warmup}[4]
  \item Oxidation occurs at the: \answerblank[2cm]{anode}
  \item Galvanic cells have $E^\circ_{\text{cell}}$:
        \answerblank[2cm]{positive}
  \item Electrons flow from \answerblank[4cm]{anode to cathode}
  \item $\Delta G^\circ$ for a galvanic cell is:
        \answerblank[2cm]{negative}
\end{warmup}

\segment{INSTRUCTION A}{25}{A scale built on an arbitrary zero}

\notesection{Standard reduction potentials}{9.9}
\practice{5}

Only a \emph{difference} in potential is measurable, so one half-reaction is
assigned a value by convention:
\[ \ce{2H+(aq) + 2e^- -> H2(g)} \qquad E^\circ \equiv \mathbf{0.00}~\text{V} \]
Every other half-reaction is measured against this
\term{standard hydrogen electrode}. \textbf{Standard conditions} means all
solutes at \SI{1}{\Molar}, gases at \SI{1}{\atm}, and \SI{25}{\celsius}.

\begin{apctrap}
\textbf{Every table is written as reductions.} Reversing a half-reaction
\answerblank[2.8cm]{flips the sign} of $E^\circ$. But multiplying a
half-reaction by a coefficient does \textbf{not} change $E^\circ$ at all ---
potential is an \term{intensive} property, volts per unit charge, not volts
per mole. Scaling a half-reaction to balance electrons is required; scaling
its voltage is wrong.
\end{apctrap}

\[ \boxed{\;E^\circ_{\text{cell}} = E^\circ_{\text{cathode}}
   - E^\circ_{\text{anode}}\;} \]
with both values taken straight from the table as reductions. For a
galvanic cell you want $E^\circ_{\text{cell}} > 0$, so the species with the
\answerblank[3cm]{more positive} $E^\circ$ is the one reduced.

\begin{workedexample}[Worked example 7: the Daniell cell]
Combine \ce{Zn^2+}/\ce{Zn} ($-0.76$~V) and \ce{Cu^2+}/\ce{Cu} ($+0.34$~V).
Copper is more positive, so copper is reduced (cathode) and zinc oxidized
(anode):
\[ \ce{Zn + Cu^2+ -> Zn^2+ + Cu} \qquad
   E^\circ_{\text{cell}} = 0.34 - (-0.76) = \mathbf{+1.10}~\text{V} \]
\end{workedexample}

\begin{workedexample}[Worked example 8: when the electrons do not balance]
Combine \ce{Ag+}/\ce{Ag} ($+0.80$~V) and \ce{Cu^2+}/\ce{Cu} ($+0.34$~V).
Silver is reduced; balancing electrons needs \emph{two} silver
half-reactions:
\[ \ce{Cu + 2Ag+ -> Cu^2+ + 2Ag} \qquad
   E^\circ_{\text{cell}} = 0.80 - 0.34 = \mathbf{+0.46}~\text{V} \]
The silver half-reaction was doubled; its $E^\circ$ was not. Writing
$2(0.80) - 0.34 = 1.26$~V is the classic error here.
\end{workedexample}

\segment{GUIDED PRACTICE}{15}{Building cells from the table}

$E^\circ$ (V): \ce{Ag+}/\ce{Ag} $+0.80$ \textbullet\ \ce{Cu^2+}/\ce{Cu}
$+0.34$ \textbullet\ \ce{Pb^2+}/\ce{Pb} $-0.13$ \textbullet\
\ce{Ni^2+}/\ce{Ni} $-0.23$ \textbullet\ \ce{Fe^2+}/\ce{Fe} $-0.44$
\textbullet\ \ce{Zn^2+}/\ce{Zn} $-0.76$ \textbullet\ \ce{Mg^2+}/\ce{Mg}
$-2.37$

\begin{enumerate}[leftmargin=*,itemsep=0.35em]
  \item \ce{Zn} and \ce{Ag+}: \answerblank[2.6cm]{$+1.56$ V}
  \item \ce{Mg} and \ce{Cu^2+}: \answerblank[2.6cm]{$+2.71$ V}
  \item \ce{Zn} and \ce{Ni^2+}: \answerblank[2.6cm]{$+0.53$ V}
  \item \ce{Fe} and \ce{Cu^2+}: \answerblank[2.6cm]{$+0.78$ V}
  \item \ce{Pb} and \ce{Ag+}: \answerblank[2.6cm]{$+0.93$ V}
\end{enumerate}

\segment{INSTRUCTION B}{20}{Volts into kilojoules, and on to $K$}

\notesection{$\Delta G^\circ = -nFE^\circ$}{9.9}
\practice{5}

Electrical work is charge moved through a potential difference. With
\term{Faraday's constant} $F = \SI{96485}{\coulomb\per\mole}$,
\[ \boxed{\;\Delta G^\circ = -nFE^\circ_{\text{cell}}\;} \]

\begin{apctrap}
\textbf{Here is where $n$ finally matters.} Voltage does not depend on how
many electrons are transferred; free energy does, because more electrons
through the same potential means more energy. Two cells can share an
$E^\circ$ and have very different $\Delta G^\circ$.

And $n$ comes from the \textbf{balanced overall equation}, not from either
half-reaction alone.
\end{apctrap}

\begin{workedexample}[Worked example 9: the Daniell cell in kilojoules]
$E^\circ = +1.10$~V and $n = 2$:
\[ \Delta G^\circ = -(2)(96485)(1.10) = -212{,}267~\text{J}
   = \mathbf{-212}~\text{kJ} \]
The answer arrives in \textbf{joules} --- $F$ is coulombs per mole and a
volt is a joule per coulomb --- so convert to kJ.
\end{workedexample}

Combining with $\Delta G^\circ = -RT\ln K$ gives the third link:
\[ nFE^\circ = RT\ln K \qquad\Longrightarrow\qquad
   \ln K = \frac{nFE^\circ}{RT} \]

\begin{center}
\renewcommand{\arraystretch}{1.2}
\begin{tabular}{@{}cccl@{}}
\hline
$E^\circ_{\text{cell}}$ & $\Delta G^\circ$ & $K$ & \\
\hline
$> 0$ & $< 0$ & $> 1$ & favored; galvanic \\
$= 0$ & $= 0$ & $= 1$ & at equilibrium; a dead battery \\
$< 0$ & $> 0$ & $< 1$ & unfavored; electrolytic \\
\hline
\end{tabular}
\end{center}

\segment{APPLICATION}{20}{Moving between the three}

\begin{enumerate}[leftmargin=*,itemsep=0.4em]
  \item For \ce{Cu + 2Ag+ -> Cu^2+ + 2Ag}, $E^\circ = +0.46$~V. Calculate
        $\Delta G^\circ$. \workspace[2.2cm]
        \keyonly{\begin{apcanswer}$n = 2$;
        $\Delta G^\circ = -(2)(96485)(0.46) = -88{,}766~\text{J}
        = \mathbf{-89}~\text{kJ}$\end{apcanswer}}
  \item Two cells both have $E^\circ = +0.50$~V; one transfers 1 electron
        and the other 4. Compare their voltages and their $\Delta G^\circ$.
        \answerlines[3]{The voltages are equal --- potential is intensive
        and independent of $n$. The free energy changes are not: the
        four-electron cell has four times the magnitude of
        $\Delta G^\circ$, since four times as much charge moves through the
        same potential.}
  \item A battery is described as ``dead''. Give $E_{\text{cell}}$,
        $\Delta G$, and the relationship between $Q$ and $K$.
        \answerblank[7.4cm]{$E = 0$, $\Delta G = 0$, $Q = K$}
\end{enumerate}

\begin{exitticket}
A student doubles a half-reaction to balance electrons and doubles its
$E^\circ$ too. Explain why the second step is wrong.
\keyonly{\begin{apcanswer}Potential is intensive --- energy per unit charge
--- so it does not scale with amount. Doubling the half-reaction doubles
both the energy and the charge, leaving the volts unchanged. Only $n$
changes, and that enters later through
$\Delta G^\circ = -nFE^\circ$.\end{apcanswer}}
\end{exitticket}

% =====================================================================
\blocklesson{Nonstandard Conditions and Electrolysis}{CED 9.10, 9.11 \textbullet\ Zumdahl \S18.4, \S18.7}

\begin{objectives}
  \item Predict qualitatively how concentration changes a cell potential.
  \item Apply Faraday's law to an electrolysis calculation.
\end{objectives}

\reading{Zumdahl \S18.4, \S18.7}{895--900, 906--911}

\begin{warmup}[4]
  \item $\Delta G^\circ = $ \answerblank[3cm]{$-nFE^\circ$}
  \item $F = $ \answerblank[3.4cm]{\SI{96485}{\coulomb\per\mole}}
  \item $Q > K$ means the reaction shifts:
        \answerblank[2.4cm]{reverse}
  \item $E^\circ_{\text{cell}}$ for a dead battery:
        \answerblank[1.4cm]{0}
\end{warmup}

\segment{INSTRUCTION A}{25}{Voltage as distance from equilibrium}

\notesection{Changing the concentrations}{9.10}
\practice{6}

A cell's voltage measures how far the reaction is from equilibrium. As it
runs, reactants are consumed and products build up, so $Q$
\answerblank[1.6cm]{rises} toward $K$ and the voltage
\answerblank[1.6cm]{falls}. When $Q = K$ the cell is dead.

\begin{apcnote}[What CED 9.10 does and does not require]
Topic 9.10 lists \textbf{no equation}. The Nernst equation is not required,
and it does not appear on the AP equation sheet. What is required is the
\emph{direction} of the effect, argued from Le Ch\^atelier:

\begin{itemize}[leftmargin=*,itemsep=0.2em]
  \item Increase a \textbf{reactant} concentration $\Rightarrow$ $Q$
        decreases $\Rightarrow$ the cell is further from equilibrium
        $\Rightarrow$ voltage \answerblank[1.8cm]{increases}.
  \item Increase a \textbf{product} concentration $\Rightarrow$ $Q$
        increases $\Rightarrow$ voltage \answerblank[1.8cm]{decreases}.
\end{itemize}

Reason with $Q$ against $K$. That argument earns full credit and needs no
logarithms.
\end{apcnote}

A \term{concentration cell} has the \emph{same} species in both half-cells
at different concentrations. $E^\circ_{\text{cell}} = $
\answerblank[1.4cm]{0}, since both electrodes are identical, yet the cell
still produces a voltage --- the system moves to equalize the two
concentrations. The \answerblank[3.4cm]{dilute} half-cell is the anode
(metal dissolves, raising its concentration) and the concentrated one is
the cathode. The cell dies when the concentrations become equal.

\segment{INSTRUCTION B}{20}{Counting electrons with an ammeter}

\notesection{Faraday's law}{9.11}
\practice{5}

Electrolysis is stoichiometry in which one reagent is
\answerblank[2.4cm]{the electron}. Current is how you count them:
\[ I = \frac{q}{t} \qquad\text{so}\qquad q = It \]
with $q$ in coulombs, $I$ in amperes, $t$ in
\answerblank[1.8cm]{seconds}. Then
$\text{mol e}^- = q/96485$.

\begin{center}
\fbox{\parbox{0.92\linewidth}{\centering
current \& time $\to$ charge ($q = It$) $\to$ moles of electrons
($\div F$) \\
$\to$ moles of substance ($\div n$ from the half-reaction) $\to$ grams
($\times$ molar mass)}}
\end{center}

\begin{apctrap}
\textbf{The half-reaction supplies the electron ratio, and it is not always
2.} One mole of electrons deposits one mole of \ce{Ag}, but only half a
mole of \ce{Cu} and a third of a mole of \ce{Al}. Reading that number off
the ion's charge is the step students skip.

Also: \textbf{time must be in seconds.} A problem stated in minutes or
hours is stated that way on purpose.
\end{apctrap}

\begin{workedexample}[Worked example 10: silver plating]
\SI{2.00}{\ampere} for \SI{30.0}{\minute} through \ce{AgNO3}.

$t = 30.0 \times 60 = \SI{1800}{\second}$;
$q = (2.00)(1800) = \SI{3600}{\coulomb}$;
$\text{mol e}^- = 3600/96485 = 0.0373$.

$\ce{Ag+ + e^- -> Ag}$ is one electron per atom, so
$m = (0.0373)(107.87) = \mathbf{4.03}$~g.
\end{workedexample}

\begin{workedexample}[Worked example 11: the same current, a different ion]
\SI{1.50}{\ampere} for \SI{1.00}{\hour} through \ce{CuSO4}.

$q = (1.50)(3600) = \SI{5400}{\coulomb}$;
$\text{mol e}^- = 5400/96485 = 0.0560$.

But $\ce{Cu^2+ + 2e^- -> Cu}$ needs \emph{two} electrons per atom:
\[ \text{mol Cu} = \frac{0.0560}{2} = 0.0280 \qquad
   m = (0.0280)(63.55) = \mathbf{1.78}~\text{g} \]
Forgetting to divide by 2 gives 3.56 g --- exactly double, and a lost
point.
\end{workedexample}

\segment{APPLICATION}{20}{Predicting and scaling}

\begin{enumerate}[leftmargin=*,itemsep=0.4em]
  \item For \ce{Zn + Cu^2+ -> Zn^2+ + Cu}, predict the effect on
        $E_{\text{cell}}$ of increasing $[\ce{Cu^2+}]$, and justify without
        calculating. \answerlines[3]{\ce{Cu^2+} is a reactant, so raising
        its concentration lowers $Q$ and moves the system further from
        equilibrium. The driving force increases, so $E_{\text{cell}}$
        increases.}
  \item How long, in hours, to deposit \SI{5.00}{\gram} of copper at
        \SI{2.00}{\ampere}? \workspace[2.6cm]
        \keyonly{\begin{apcanswer}mol Cu $= 5.00/63.55 = 0.0787$;
        mol e$^- = 0.157$; $q = (0.157)(96485) = \SI{15183}{\coulomb}$;
        $t = 15183/2.00 = \SI{7591}{\second} =
        \mathbf{2.11}~\text{h}$\end{apcanswer}}
  \item A concentration cell is built from two \ce{Cu} electrodes in
        \SI{0.10}{\Molar} and \SI{1.0}{\Molar} \ce{Cu^2+}. Identify the
        anode and say when the cell stops.
        \answerlines[3]{The \SI{0.10}{\Molar} half-cell is the anode:
        copper dissolves there, raising its concentration, while
        \ce{Cu^2+} plates out in the concentrated half-cell. The cell stops
        when the two concentrations become equal, at which point $Q = K$
        and $E_{\text{cell}} = 0$.}
\end{enumerate}

\begin{exitticket}
State the one question to ask before starting any calculation in this unit,
and why it settles the method.
\keyonly{\begin{apcanswer}``Which of $\Delta G^\circ$, $K$, and
$E^\circ_{\text{cell}}$ am I given, and which is wanted?'' All three
describe the same fact and are linked by
$\Delta G^\circ = -RT\ln K = -nFE^\circ$, so naming the two endpoints
selects the equation and leaves only unit conversion to get
wrong.\end{apcanswer}}
\end{exitticket}

\end{document}
